\chapter{Discussions}
\label{chap:discussions}

Ce chapitre interprète les résultats présentés au chapitre~\ref{chap:resultats} à la lumière
des deux hypothèses de l'étude et de leur portée réelle. La section~\ref{sec:disc-h1} traite
de la signification du taux de non-conformité observé ; la section~\ref{sec:disc-h2} discute
la robustesse du module et la valeur informative d'un intervalle de confiance de largeur nulle ;
la section~\ref{sec:disc-hors-domaine} délimite le périmètre d'usage licite de la fonction ;
les sections~\ref{sec:disc-validite-interne} et~\ref{sec:disc-validite-externe} examinent
respectivement les menaces à la validité interne et les conditions de généralisation.

% ------------------------------------------------------------------
\section{L'artefact respecte-t-il sa spécification ? (H1)}
\label{sec:disc-h1}

\subsection{Interprétation du taux de non-conformité}
\label{subsec:disc-h1-taux}

Sur 12\,000 essais, 200 présentent au moins un écart entre le moteur JavaScript
\texttt{computeRFM} et la référence Python calculant les sommes en centimes entiers, soit un
taux de non-conformité de $1{,}667\,\%$ (IC Monte-Carlo à $95\,\%$ : $[1{,}448\,\% ;\, 1{,}904\,\%]$,
tableau~\ref{tab:resultats-conformite-bilan}). L'intervalle confirme que cette valeur n'est
pas un artefact de l'échantillonnage Monte-Carlo : la borne supérieure reste inférieure à
$2\,\%$, seuil opérationnel implicitement retenu dans le protocole expérimental comme limite
de tolérance pour le contexte de petits commerçants congolais.

La question est double : ce taux est-il statistiquement significatif, et est-il acceptable
pour l'usage auquel le module est destiné ? Sur le premier point, la réponse est affirmative ;
200 essais non conformes sur 12\,000 ne relèvent pas du bruit d'estimation et signalent un
mécanisme systématique reproductible. Sur le second point, l'appréciation dépend du contexte
opérationnel. Un taux de conformité de $98{,}3\,\%$ demeure élevé pour un module de
segmentation destiné à des PME informelles dont les données d'entrée sont elles-mêmes soumises
à des imprécisions de saisie. La segmentation RFM sert ici à orienter des actions commerciales
— identifier les clients à fidéliser, ceux à réactiver — et non à produire des états
comptables certifiés. Dans ce cadre décisionnel, une erreur de segmentation portant sur le
score $M$ d'un client isolé, dans une situation de montants proches d'un seuil de tercile,
est susceptible de n'avoir aucune conséquence pratique sur la stratégie commerciale adoptée.

Il convient néanmoins de ne pas minimiser ce résultat : les erreurs observées sont
déterministes et reproductibles à partir des mêmes graines, et les $1{,}667\,\%$ traduisent
une fragilité structurelle du moteur, non un bruit aléatoire corrigeable par répétition.

\subsection{Cause identifiée : accumulation en virgule flottante}
\label{subsec:disc-h1-cause}

Les 489 occurrences de score $M$ erroné et les 2 occurrences d'étiquette erronée qui en
découlent ont une cause unique et identifiable : l'addition des montants de commandes est
effectuée en virgule flottante IEEE~754 double précision dans le moteur TypeScript, alors que
la référence Python additionne les mêmes montants convertis en centimes entiers. Lorsque
plusieurs commandes sont agrégées, l'erreur d'arrondi cumulée peut modifier la relation
d'ordre entre deux clients dont les sommes affichées sont identiques. Le contre-exemple
minimal est reproduit au tableau~\ref{tab:resultats-contre-exemple} : les clients 00001 et
00002 affichent tous deux $30{,}05$~USD après trois commandes, mais leur représentation
interne en virgule flottante diffère d'un epsilon infime, ce qui déplace le score $M$ du
client~00002 de~1 à~2.

Ce type d'erreur est bien documenté dans la littérature sur l'arithmétique en virgule flottante
\citep{Haissinsky2012} et dans les guides de pratique numérique \citep{Ribatet_nd}. Il n'est
pas propre à TypeScript ou à JavaScript ; il affecte tout langage reposant sur la norme
IEEE~754 sans couche d'arrondi explicite au niveau des sommes. Des bibliothèques de calcul
monétaire comme \texttt{decimal.js} ou \texttt{dinero.js} apportent une réponse classique,
au prix d'une dépendance supplémentaire et d'une migration des structures de données internes.
La correction pourrait également être réalisée nativement, en arrondissant chaque somme
partielle au centime le plus proche avant toute comparaison de quantiles, ce qui supprimerait
le mécanisme identifié sans modifier l'algorithme de segmentation.

\paragraph{Ce que H1 prouve et ce qu'il ne prouve pas.}
La conformité mesurée dans cette étude porte sur la logique de segmentation RFM de
\texttt{computeRFM} au commit \texttt{b062a93a}, comparée à une référence indépendante en
Python. Elle démontre que le moteur reproduit fidèlement les règles de tercilisation dans
$98{,}3\,\%$ des configurations testées et identifie une cause précise dans les $1{,}7\,\%$
restants. Elle ne constitue pas une évaluation de l'utilité commerciale des segments produits,
de leur pertinence prédictive pour les comportements d'achat réels, ni de leur adéquation au
marché congolais. Ces questions relèvent d'une validation de contenu \citep{Salminen2023} et
d'une validation prospective \citep{Fader2005} qui excèdent le périmètre de la présente
campagne.

Il convient également de distinguer la conformité de l'implémentation de la pertinence du
choix de la méthode RFM elle-même. La littérature récente \citep{ViannaFilho2026,
Kassambara2026} soulève régulièrement la question de savoir si une segmentation en terciles
sur trois dimensions constitue la meilleure approche pour capturer la hétérogénéité des
comportements clients. Ces débats sont légitimes, mais ils portent sur la conception de
l'artefact — en particulier sur le choix des seuils et du nombre de groupes — et non sur
la conformité de l'implémentation à sa propre spécification. La présente évaluation ne
prend pas position sur ce premier niveau de choix ; elle évalue uniquement la fidélité de
l'exécution à la règle telle qu'elle est définie dans la codebase d'ELMES-TEKA.

% ------------------------------------------------------------------
\section{Le module est-il robuste aux imprécisions courantes ? (H2)}
\label{sec:disc-h2}

\subsection{Un résultat plus fort qu'attendu}
\label{subsec:disc-h2-force}

Les quatre cellules admissibles pour le critère H2 — tailles $n \in \{50,\, 200\}$,
profils \texttt{regular} et \texttt{skewed}, bruit monétaire $\pm1\,\%$, clients éloignés
des seuils selon la règle à $5\,\%$ — présentent toutes un taux de changement moyen nul,
avec des intervalles de confiance ponctuels $[0\,\% ;\, 0\,\%]$
(tableau~\ref{tab:resultats-h2}). La borne supérieure de chaque intervalle est strictement
inférieure au seuil de $5\,\%$ fixé par le protocole ; la décision est «~Compatible~» pour
les quatre cellules.

Ce résultat est plus fort que ce qui était strictement nécessaire pour satisfaire H2. Le
critère aurait été respecté avec des taux inférieurs à $5\,\%$ ; l'observation de $0\,\%$
exact dans toutes les répétitions n'était pas attendue a priori et mérite une explication.

\subsection{Pourquoi les clients éloignés sont stables}
\label{subsec:disc-h2-pourquoi}

La stabilité complète s'explique par la définition même de la condition d'éloignement. Un
client est dit éloigné si sa distance normalisée à chacun des six seuils initiaux ($s_{R1}$,
$s_{R2}$, $s_{F1}$, $s_{F2}$, $s_{M1}$, $s_{M2}$) est supérieure à $5\,\%$. Une
perturbation monétaire de $\pm1\,\%$ ne peut pas déplacer un seuil de tercile au-delà de
$5\,\%$ sur la dimension $M$ lorsque les clients sont distribués de façon régulière ou
dissymétrique, car le déplacement d'un quantile dépend de l'ensemble de la distribution et
non d'un seul client. Un client situé à plus de $5\,\%$ de chaque seuil possède donc une
«~marge de sécurité~» qui absorbe le bruit de $\pm1\,\%$ sans franchir le seuil le plus
proche. Dans ce sens, la condition d'éloignement est une garantie théorique de stabilité,
et le résultat observé confirme empiriquement que la campagne n'a pas rencontré de cas où
cette garantie aurait été mise en défaut dans les 100 répétitions par cellule.

\paragraph{La figure~\ref{fig:ann-sensibilite-montants} (Annexe~\ref{ann:reproduction})
illustre ce mécanisme} : les taux de changement global augmentent avec l'intensité du bruit
monétaire pour les profils \texttt{regular} et \texttt{skewed}, mais les clients éloignés
ne changent pas sous $\pm1\,\%$.

\subsection{Signification d'un intervalle de confiance $[0\,\% ;\, 0\,\%]$}
\label{subsec:disc-h2-ic}

L'intervalle de confiance $[0\,\% ;\, 0\,\%]$ signifie que, parmi les 100 répétitions de
chaque cellule, aucune n'a produit un changement de segment pour un client éloigné. Il
traduit une absence de variation observée, non une démonstration qu'aucun changement n'est
possible dans n'importe quelle configuration d'usage. Comme rappelé à la
section~\ref{sec:resultats-cibles}, un intervalle Monte-Carlo de largeur nulle ne constitue
pas une preuve d'une probabilité de changement identiquement nulle ; il signifie simplement
que l'estimateur ponctuel et ses deux bornes coïncident au zéro empirique.

Cette précision est importante pour éviter une surinterprétation : le résultat ne garantit
pas la robustesse du module pour des corpus réels très différents des mécanismes générateurs
synthétiques utilisés dans la campagne, ni pour des niveaux de bruit supérieurs à $\pm1\,\%$.
La figure~\ref{fig:ann-sensibilite-montants} montre que pour les profils non-ties, les taux
de changement global augmentent graduellement de $\pm1\,\%$ à $\pm20\,\%$ de bruit monétaire,
ce qui confirme que la robustesse est conditionnelle au niveau de perturbation.

La décision de construire l'intervalle par loi $t$ plutôt que par bootstrap s'explique par
la symétrie des estimateurs et le nombre de répétitions ($n = 100$ par cellule), suffisant
pour invoquer le théorème central limite \citep{Diez2019}. Sur une variable dont la variance
empirique est nulle, les deux méthodes convergent vers le même résultat ; la distinction est
sans conséquence pratique pour les cellules H2, mais importante à rappeler pour des cellules
où la variance n'est pas nulle. Les données cellule par cellule, disponibles dans l'archive
\texttt{hypothesis\_H2.csv} (Annexe~\ref{ann:reproduction}), permettraient une vérification
par bootstrap si une révision du protocole le jugeit utile.

\subsection{Positionnement dans la littérature}
\label{subsec:disc-h2-litterature}

La question de la robustesse d'une segmentation RFM aux imprécisions des données d'entrée est
peu traitée comme objet explicite d'étude dans la littérature récente. \citet{Rajendran2024}
et \citet{Sobantu2025} proposent des évaluations de méthodes de segmentation RFM dans des
contextes applicatifs variés, mais sans critère formel équivalent au taux de changement sous
perturbation contrôlée. \citet{Mingmongkhol2023} analyse la sensibilité d'une segmentation
RFM à la pondération des dimensions, sans soumettre l'algorithme à des erreurs de données
synthétiques. Cette étude adopte une approche différente, plus proche des protocoles de
simulation de robustesse utilisés en évaluation de méthodes statistiques \citep{Morris2019},
ce qui rend la comparaison directe difficile mais justifie l'originalité de la démarche
adoptée.

% ------------------------------------------------------------------
\section{Quelles situations sortent du domaine d'usage ?}
\label{sec:disc-hors-domaine}

\subsection{Le profil identique (\emph{ex~æquo}) : un signal de qualité des données}
\label{subsec:disc-ties}

Le profil \texttt{ties} — tous les clients présentant exactement le même nombre de commandes,
la même date de dernier achat et le même ticket unitaire — n'est pas inclus dans le critère
H2. Cette exclusion est une décision du protocole expérimental, motivée avant l'analyse des
résultats par le fait que ce profil se situe hors du domaine d'usage normal de la fonction
\texttt{computeRFM}. Elle ne résulte pas d'une observation \emph{a posteriori} visant à
retirer des données gênantes.

Le comportement observé sur ce profil justifie cette décision. Sous perturbation de dates,
les taux de changement atteignent 42 à 63~\% selon la taille du corpus
(tableau~\ref{tab:resultats-ties-dates}). L'explication est structurelle : lorsque tous les
clients ont les mêmes valeurs $R$, $F$ et $M$, les seuils de tercile ne sont pas définis de
façon unique — tout point de la plage commune entre deux clients constitue un seuil valide.
L'ajout d'un bruit infime sur les dates modifie les relations d'ordre entre clients de façon
arbitraire, ce qui déplace les seuils aléatoirement et produit une recomposition complète des
groupes. Ce comportement n'est pas un défaut de l'algorithme ; il reflète l'indétermination
mathématique des quantiles sur des valeurs identiques.

En pratique, un corpus réel présentant des valeurs RFM entièrement homogènes signale un
problème de qualité des données en amont : doublons de clients, saisie en lot à valeur fixe,
ou consolidation erronée. Le profil \texttt{ties} constitue ainsi un cas limite qui
n'invalide pas H2 mais délimite le périmètre à l'intérieur duquel la robustesse mesurée est
pertinente. Les cinq comportements hors domaine identifiés dans la campagne
(tableau~\ref{tab:resultats-hors-domaine}) — mélange de devises, identifiant absent, date
invalide, date future, montant négatif — confirment que la fonction \texttt{computeRFM}
n'effectue pas de contrôle métier en entrée et s'appuie sur les couches amont du SaaS pour
garantir la qualité des données transmises.

\subsection{Le corpus de taille $n = 3$ : seuil minimal de stabilité}
\label{subsec:disc-n3}

Les essais conduits sur des corpus de trois clients ont servi à la mesure de conformité (H1)
mais ont été exclus du critère H2. À cette taille, le premier tiers d'un groupe de trois
clients correspond à un unique client, ce qui rend les seuils de tercile extrêmement sensibles
à toute modification individuelle. Le calcul reste mathématiquement valide, mais la
segmentation obtenue varie fortement d'une répétition à l'autre, non en raison d'un défaut
de l'algorithme, mais parce que la notion de «~troisième~» n'est pas stable sur un ensemble
de trois éléments.

Cette limite n'est pas propre à ELMES-TEKA ; elle est inhérente à toute méthode de
segmentation par quantiles. La recommandation opérationnelle qui en découle est qu'un corpus
d'au moins une dizaine de clients constitue un minimum raisonnable pour que les segments RFM
soient interprétables de façon stable. \citet{Yoseph2023} et \citet{Ray2023} utilisent
des corpus de plusieurs centaines à plusieurs milliers de clients dans leurs études
applicatives, sans que cette contrainte minimale soit explicitement discutée ; elle semble
traitée comme une évidence pratique. La présente campagne la confirme de façon expérimentale.

% ------------------------------------------------------------------
\section{La méthode de validation est-elle appropriée ? (menaces à la validité interne)}
\label{sec:disc-validite-interne}

\subsection{Forces du protocole}
\label{subsec:disc-vi-forces}

Le protocole expérimental présente plusieurs caractéristiques qui renforcent la validité
interne des résultats. Les graines de génération et de perturbation sont explicitement
consignées pour chaque essai, ce qui garantit la reproductibilité complète de la campagne.
L'appariement systématique de chaque corpus perturbé à son corpus initial d'origine supprime
la variabilité inter-corpus qui aurait confondé les effets de la perturbation avec des
différences structurelles entre jeux de données. Le cadre \emph{ADEMP} proposé par
\citet{Morris2019} a été suivi comme structure de présentation, ce qui facilite la
comparaison avec d'autres études de simulation et rend les choix méthodologiques explicites.

La référence Python indépendante est construite sur un principe arithmétique différent du
moteur évalué — addition en centimes entiers contre addition en virgule flottante —, ce qui
lui confère une capacité de détection ciblée sur le mécanisme de non-conformité identifié.
La disponibilité de l'archive expérimentale complète avec empreinte SHA-256 \citep{ViannaFilho2026}
assure la traçabilité des résultats jusqu'à leur source.

\subsection{Limites de la référence}
\label{subsec:disc-vi-limites-ref}

La référence Python n'est pas une implémentation certifiée par une tierce partie indépendante.
Elle a été développée pour les besoins de la campagne à partir de la spécification de
\texttt{computeRFM} et pourrait contenir des erreurs propres qui ne seraient pas détectées
par comparaison avec le moteur TypeScript. Les neuf cas manuels (section~\ref{subsec:resultats-cas-manuels})
ont été vérifiés contre la spécification et donnent confiance dans le comportement de la
référence sur les configurations limites, mais cette vérification ne couvre pas l'ensemble de
l'espace des entrées.

Cette limite est structurelle dans toute démarche de validation par référence construite à
la volée ; elle appelle à terme une certification de la référence par un tiers ou une
vérification formelle de son algorithme de tercilisation.

À titre de contrôle partiel, les neuf cas manuels ont été vérifiés indépendamment de la
référence Python, contre la spécification textuelle de l'algorithme. Leur accord à $9/9$
confirme que la référence ne commet pas d'erreur sur les configurations limites les plus
critiques — corpus vide, client unique, valeurs toutes identiques, filtre d'exercice, montant
nul. Pour autant, cette vérification ne couvre pas les configurations de la pleine campagne,
au sein desquelles la référence est utilisée comme oracle sans contrôle additionnel.

\subsection{Profils synthétiques non calibrés}
\label{subsec:disc-vi-synth}

Les profils de génération — \texttt{regular}, \texttt{skewed}, \texttt{ties} — ont été
construits pour couvrir des distributions contrastées, mais ils n'ont pas été calibrés sur
des données de transactions réelles de petits commerçants congolais. Les valeurs de montants,
les distributions de fréquence et les horizons temporels utilisés dans la campagne peuvent
s'éloigner sensiblement de ce qui est observé dans un portefeuille réel d'ELMES-TEKA. Les
résultats décrivent donc le comportement de la fonction sur ces mécanismes générateurs
spécifiques, non son comportement en production, comme le souligne la contrainte méthodologique
rappelée à la section~\ref{sec:resultats-hors-domaine}.

\subsection{Périmètre de l'évaluation}
\label{subsec:disc-vi-perimetre}

La campagne évalue la fonction \texttt{computeRFM} de façon isolée, à l'interface entre des
données de transactions en entrée et des vecteurs de scores en sortie. Elle ne mesure pas le
comportement de la couche MongoDB, des routes HTTP, du système d'authentification multi-tenant,
ni des effets économiques de la segmentation sur les commerçants. Ces dimensions font partie
de l'artefact complet ELMES-TEKA \citep{James2023} mais dépassent le périmètre défini par
les hypothèses H1 et H2.

% ------------------------------------------------------------------
\section{Quelles généralisations sont licites ? (menaces externes et perspectives)}
\label{sec:disc-validite-externe}

\subsection{Limites de généralisation}
\label{subsec:disc-ve-limites}

Les résultats sont valables pour la logique de \texttt{computeRFM} telle que codée au commit
\texttt{b062a93a7be96d082dee9afe66b9f4f1f9ebecaa}. Toute modification du code de segmentation
— changement du nombre de groupes, remplacement des terciles par d'autres quantiles, ajout
d'une couche d'arrondi monétaire — invaliderait les mesures obtenues et nécessiterait une
nouvelle campagne de validation. Ce point est d'importance pratique : si le correctif
arithmétique recommandé par H1 est appliqué, il conviendrait de confirmer que le taux de
non-conformité est bien ramené à zéro dans les conditions ayant produit des écarts.

Les données synthétiques ne sont pas représentatives de l'ensemble des commerçants congolais
utilisant ELMES-TEKA. Les distributions de montants, les fréquences d'achat et les profils
de récence varient selon le secteur d'activité, la taille du marché et les habitudes de
paiement locales. Une conclusion du type «~le module est robuste pour les petits commerçants
congolais~» irait au-delà de ce que les données synthétiques permettent d'établir.

\subsection{Perspectives de validation}
\label{subsec:disc-ve-perspectives}

Plusieurs directions permettraient d'étendre la portée de la présente évaluation. Une
validation sur données réelles anonymisées, extraites de la base de production d'ELMES-TEKA,
permettrait de calibrer les mécanismes générateurs et de mesurer la sensibilité du module
sur des distributions effectives \citep{Salminen2023}. Cette validation soulèverait
toutefois des contraintes de confidentialité liées à la nature multi-tenant de la plateforme.

L'extension à un contexte multi-devises — en particulier la cohabitation du franc congolais
(CDF) et du dollar américain (USD), pratique courante dans le commerce kinois — constituerait
une priorité métier. Le comportement hors domaine observé lors du mélange de devises
(tableau~\ref{tab:resultats-hors-domaine}) montre que la fonction traite actuellement toutes
les valeurs comme commensurables, ce qui produit des montants agrégés économiquement
incohérents.

Une validation prospective comparerait les segments produits par \texttt{computeRFM} avec
les comportements de réachat ultérieurs des clients, testant ainsi la validité prédictive des
étiquettes \citep{Fader2005}. Ce type de validation excède la portée d'une évaluation de
conformité algorithmique, mais il est indispensable pour établir la valeur décisionnelle des
segments dans le contexte opérationnel. Des modèles stochastiques de valeur client tels que
Pareto/NBD \citep{Fader2005} ou les variantes bayésiennes récentes pourraient fournir une
référence externe pour juger de la qualité discriminante des cinq catégories retenues.

Enfin, l'incorporation d'analyses de sensibilité globale \citep{Lakens2013} permettrait
d'estimer la contribution relative des dimensions $R$, $F$ et $M$ à la variabilité
inter-sessions de la segmentation, au-delà de l'analyse par scénario conduite dans la
présente campagne. Ces pistes constituent des prolongements naturels, conditionnels à la
disponibilité de données réelles et à l'accord des parties prenantes d'ELMES-TEKA sur les
modalités d'accès.

\subsection{Récapitulatif des conclusions par hypothèse}
\label{subsec:disc-recap}

Le tableau~\ref{tab:disc-recap} synthétise les décisions retenues pour H1 et H2 à l'issue
de la discussion.

\begin{table}[htb]
  \centering
  \caption{Récapitulatif des décisions d'hypothèse après discussion.}
  \label{tab:disc-recap}
  \begin{tabular}{lp{0.52\textwidth}l}
    \toprule
    Hypothèse & Résultat et décision & Condition \\
    \midrule
    H1 & Taux de non-conformité $1{,}667\,\%$ (IC $[1{,}448\,\% ;\, 1{,}904\,\%]$),
         inférieur au seuil de $2\,\%$ ; cause identifiée (virgule flottante) ;
         décision : \textbf{conforme sous réserve}
      & Commit \texttt{b062a93a} \\
    H2 & Taux de changement $0{,}000\,\%$ dans les 4 cellules admissibles (IC $[0\,\% ;\, 0\,\%]$) ;
         borne supérieure $< 5\,\%$ ; décision : \textbf{compatible}
      & Clients éloignés, $\pm1\,\%$, $n \ge 50$ \\
    \bottomrule
  \end{tabular}
\end{table}

La mention «~conforme sous réserve~» pour H1 signifie que la conformité est établie dans les
limites du périmètre testé et sous la cause identifiée, qui est corrigible. La décision ne
peut être portée à «~conforme sans réserve~» qu'après application du correctif arithmétique
et réévaluation de la campagne sur les essais ayant produit des écarts.
